\section{Metodologia}

\subsection{Desenho Experimental}

Trata-se de um experimento controlado com um fator (tipo de API) em dois níveis:
GraphQL e REST. As variáveis são:

\begin{itemize}
    \item \textbf{Independente:} tipo de API (categórica binária: GraphQL ou REST)
    \item \textbf{Dependentes:} tempo de resposta em milissegundos (float),
          tamanho do payload em bytes (int)
    \item \textbf{Controladas:} mesmo servidor, mesma query semântica,
          mesmo dataset, mesma máquina, mesma rede, mesmo horário de execução
\end{itemize}

\subsection{Objeto Experimental}

Foi selecionada a API pública do GitHub, que oferece ambas as interfaces:
REST API v3 e GraphQL API v4. A query consiste na obtenção de dados do
repositório \texttt{facebook/react}: nome, estrelas e forks.

\subsection{Procedimento}

Para cada trial: (1) registrar timestamp, (2) disparar requisição,
(3) medir tempo com \texttt{time.perf\_counter()}, (4) registrar tamanho do
payload com \texttt{len(response.content)}. Os trials são interleaved
(REST, GraphQL, REST, GraphQL, \dots) com intervalo de 0,5s entre eles
para mitigar deriva temporal e rate limiting.

\subsection{Critérios de Exclusão}

Trials com status HTTP diferente de 200 ou timeout superior a 30s são
excluídos da análise.

\subsection{Ameaças à Validade}

\begin{itemize}
    \item \textbf{Interna:} variação de rede, cache do servidor, garbage
          collector do Python
    \item \textbf{Externa:} resultados podem não generalizar para outras APIs
    \item \textbf{Construct:} "tamanho da resposta" mede o payload JSON bruto,
          não o volume semântico de informação transmitida
\end{itemize}
